\ifdefined\gkver\else\def\gkver{student}\fi
\documentclass[\gkver]{gaokaozhenti}
\begin{document}

\chapter{2007年广东}

\begin{enumerate}
\item
%% number: 1
%% paperName: 2007年广东高考第 1 题 4 分
%% typeId: 2
%% type: 多选题
%% chapter: 其他
%% point: 物理学史
%% method: 无
%% score: 4
%% degree: 750
%% duplicateId: 0
%% body:
许多科学家在物理学发展过程中做出了重要贡献，下列表述正确的是\xzanswer{ABD}

\fourchoices[answer=ABD]
{卡文迪许测出引力常数}
{法拉第发现电磁感应现象}
{安培提出了磁场对运动电荷的作用力公式}
{库仑总结并确认了真空中两个静止点电荷之间的相互作用规律}

%% answer: ABD
%% memo:
\memoanswer{
本题原卷及材料中未提供详解，待补充。
}

\item
%% number: 2
%% paperName: 2007年广东高考第 2 题 4 分
%% typeId: 1
%% type: 单选题
%% chapter: 原子和原子核
%% point: 玻尔模型
%% method: 无
%% score: 4
%% degree: 750
%% duplicateId: 0
%% body:
如图所示为氢原子的四个能级，其中为基态，若氢原子 A 处于激发态 $E_{2}$，氢原子 B 处于激发态 $E_{3}$，则下列说法正确的是\xzanswer{B}

\onepicture[w=3.8cm]{figs/02.png}

\fourchoices[answer=B]
{原子 A 可能辐射出 3 种频率的光子}
{原子 B 可能辐射出 3 种频率的光子}
{原子 A 能够吸收原子 B 发出的光子并跃迁道能级 $E_{4}$}
{原子 B 能够吸收原子 A 发出的光子并跃迁道能级 $E_{4}$}

%% answer: B
%% memo:
\memoanswer{
本题原卷及材料中未提供详解，待补充。
}

\item
%% number: 3
%% paperName: 2007年广东高考第 3 题 4 分
%% typeId: 2
%% type: 多选题
%% chapter: 电场
%% point: 等势面
%% method: 无
%% score: 4
%% degree: 550
%% duplicateId: 0
%% body:
如图所示的匀强电场 $E$ 的区域内，由 A、B、C、D、$A'$、$B'$、$C'$、$D'$ 作为顶点构成一正方体空间，电场方向与面 ABCD 垂直。下列说法正确的是\xzanswer{BD}

\onepicture[w=4.6cm]{figs/03.png}

\fourchoices[answer=BD]
{AD 两点间电势差 $U_{AD}$ 与 $AA'$ 两点间电势差 $U_{AA'}$ 相等}
{带正电的粒子从 A 点沿路径 $A\to D\to D'$ 移到 $D'$ 点，电场力做正功}
{带负电的粒子从 A 点沿路径 $A\to D\to D'$ 移到 $D'$ 点，电势能减小}
{带电的粒子从 A 点移到 $C'$ 点，沿对角线 $AC'$ 与沿路径 $A\to B\to B'\to C'$ 电场力做功相同}

%% answer: BD
%% memo:
\memoanswer{
本题原卷及材料中未提供详解，待补充。
}

\item
%% number: 4
%% paperName: 2007年广东高考第 4 题 4 分
%% typeId: 2
%% type: 多选题
%% chapter: 功和能
%% point: 交通工具启动
%% method: 无
%% score: 4
%% degree: 550
%% duplicateId: 0
%% body:
机车从静止开始沿平直轨道做匀加速运动，所受的阻力始终不变，在此过程中，下列说法正确的是\xzanswer{AD}

\fourchoices[answer=AD]
{机车输出功率逐渐增大}
{机车输出功率不变}
{在任意两相等时间内，机车动能变化相等}
{在任意两相等时间内，机车动量变化大小相等}

%% answer: AD
%% memo:
\memoanswer{
本题原卷及材料中未提供详解，待补充。
}

\item
%% number: 5
%% paperName: 2007年广东高考第 5 题 4 分
%% typeId: 1
%% type: 单选题
%% chapter: 力物体的平衡
%% point: 多力平衡
%% method: 无
%% score: 4
%% degree: 650
%% duplicateId: 0
%% body:
如图所示，在倾角为 $\theta$ 的固定光滑斜面上，质量为 $m$ 的物体受外力 $F_{1}$ 和 $F_{2}$ 的作用，$F_{1}$ 方向水平向右，$F_{2}$ 方向竖直向上。若物体静止在斜面上，则下列关系正确的是\xzanswer{B}

\onepicture[w=6.0cm]{figs/05.png}

\fourchoices[answer=B]
{$F_{1}\sin\theta+F_{2}\cos\theta=mg\sin\theta$，$F_{2}\leq mg$}
{$F_{1}\cos\theta+F_{2}\sin\theta=mg\sin\theta$，$F_{2}\leq mg$}
{$F_{1}\sin\theta-F_{2}\cos\theta=mg\sin\theta$，$F_{2}\leq mg$}
{$F_{1}\cos\theta-F_{2}\sin\theta=mg\sin\theta$，$F_{2}\leq mg$}

%% answer: B
%% memo:
\memoanswer{
本题原卷及材料中未提供详解，待补充。
}

\item
%% number: 6
%% paperName: 2007年广东高考第 6 题 4 分
%% typeId: 1
%% type: 单选题
%% chapter: 电场
%% point: 带电粒子在交变电场中的运动
%% method: 无
%% score: 4
%% degree: 550
%% duplicateId: 0
%% body:
平行板间加如图（a）所示周期变化的电压，重力不计的带电粒子静止在平行板中央，从 $t=0$ 时刻开始将其释放，运动过程无碰板情况。图（b）中，能定性描述粒子运动的速度图象正确的是\xzanswer{A}

\onepicture[num={（a）}, labelprefix={}, w=4.8cm]{figs/06.png}

\fourchoices[answer=A, ispicture=true, h=1.9cm]
{figs/06a.png}
{figs/06b.png}
{figs/06c.png}
{figs/06d.png}

%% answer: A
%% memo:
\memoanswer{
本题原卷及材料中未提供详解，待补充。
}

\item
%% number: 7
%% paperName: 2007年广东高考第 7 题 4 分
%% typeId: 2
%% type: 多选题
%% chapter: 交流电
%% point: 变压器
%% method: 无
%% score: 4
%% degree: 550
%% duplicateId: 0
%% body:
如图是霓虹灯的供电电路，电路中的变压器可视为理想变压器，已知变压器原线圈与副线圈匝数比 $\frac{n_{1}}{n_{2}}=\frac{1}{20}$，加在原线圈的电压为 $u_{1}=311\sin100\pi t$（$\UV$），霓虹灯正常工作的电阻 $R=440\UkO$，$I_{1}$、$I_{2}$ 表示原、副线圈中的电流，下列判断正确的是\xzanswer{BD}

\onepicture[w=7.0cm]{figs/07.png}

\fourchoices[answer=BD]
{副线圈两端电压 $6220\UV$，副线圈中的电流 $14.1\UmA$}
{副线圈两端电压 $4400\UV$，副线圈中的电流 $10.0\UmA$}
{$I_{1}<I_{2}$}
{$I_{1}>I_{2}$}

%% answer: BD
%% memo:
\memoanswer{
本题原卷及材料中未提供详解，待补充。
}

\item
%% number: 8
%% paperName: 2007年广东高考第 8 题 4 分
%% typeId: 1
%% type: 单选题
%% chapter: 电路
%% point: 传感器
%% method: 无
%% score: 4
%% degree: 350
%% duplicateId: 0
%% body:
压敏电阻的阻值随所受压力的增大而减小，右位同学利用压敏电阻设计了判断小车运动状态的装置，其工作原理如图（a）所示，将压敏电阻和一块挡板固定在绝缘小车上，中间放置一个绝缘重球。小车向右做直线运动过程中，电流表示数如图（b）所示，下列判断正确的是\xzanswer{D}

\twopicture[numA={（a）}, numB={（b）}, labelprefix={}, wA=4.8cm, wB=4.0cm]
{figs/08a.png}
{figs/08b.png}

\fourchoices[answer=D]
{从 $t_{1}$ 到 $t_{2}$ 时间内，小车做匀速直线运动}
{从 $t_{1}$ 到 $t_{2}$ 时间内，小车做匀加速直线运动}
{从 $t_{2}$ 到 $t_{3}$ 时间内，小车做匀速直线运动}
{从 $t_{2}$ 到 $t_{3}$ 时间内，小车做匀加速直线运动}

%% answer: D
%% memo:
\memoanswer{
本题原卷及材料中未提供详解，待补充。
}

\item
%% number: 9
%% paperName: 2007年广东高考第 9 题 4 分
%% typeId: 1
%% type: 单选题
%% chapter: 气体性质
%% point: 理想气体状态方程
%% method: 无
%% score: 4
%% degree: 850
%% duplicateId: 0
%% body:
一定质量的理想气体，在某一平衡状态下的压强、体积和温度分别为 $p_{1}$、$V_{1}$、$T_{1}$，在另一平衡状态下的压强、体积和温度分别为 $p_{2}$、$V_{2}$、$T_{2}$，下列关系正确的是\xzanswer{D}

\fourchoices[answer=D]
{$p_{1}=p_{2}$，$V_{1}=2V_{2}$，$T_{1}=\frac{1}{2}T_{2}$}
{$p_{1}=p_{2}$，$V_{1}=\frac{1}{2}V_{2}$，$T_{1}=2T_{2}$}
{$p_{1}=2p_{2}$，$V_{1}=2V_{2}$，$T_{1}=2T_{2}$}
{$p_{1}=2p_{2}$，$V_{1}=V_{2}$，$T_{1}=2T_{2}$}

%% answer: D
%% memo:
\memoanswer{
本题原卷及材料中未提供详解，待补充。
}

\item
%% number: 10
%% paperName: 2007年广东高考第 10 题 4 分
%% typeId: 2
%% type: 多选题
%% chapter: 气体性质
%% point: 物体的内能
%% method: 无
%% score: 4
%% degree: 750
%% duplicateId: 0
%% body:
如图为焦耳实验装置图，用绝热性能良好的材料将容器包好，重物下落带动叶片搅拌容器里的水，引起水温升高。关于这个实验，下列说法正确的是\xzanswer{AC}

\onepicture[w=4.8cm]{figs/10.png}

\fourchoices[answer=AC]
{这个装置可测定热功当量}
{做功增加了水的热量}
{做功增加了水的内能}
{功和热量是完全等价的，无区别}

%% answer: AC
%% memo:
\memoanswer{
本题原卷及材料中未提供详解，待补充。
}

\item
%% number: 11
%% paperName: 2007年广东高考第 11 题 4 分
%% typeId: 2
%% type: 多选题
%% chapter: 光学
%% point: 光的干涉
%% method: 无
%% score: 4
%% degree: 750
%% duplicateId: 0
%% body:
关于光的性质，下列说法正确的是\xzanswer{BC}

\fourchoices[answer=BC]
{光在介质中的速度大于光在真空中的速度}
{双缝干涉说明光具有波动性}
{光在同种介质种沿直线传播}
{光的偏振现象说明光是纵波}

%% answer: BC
%% memo:
\memoanswer{
本题原卷及材料中未提供详解，待补充。
}

\item
%% number: 12
%% paperName: 2007年广东高考第 12 题 4 分
%% typeId: 1
%% type: 单选题
%% chapter: 机械波
%% point: 波的图象
%% method: 无
%% score: 4
%% degree: 650
%% duplicateId: 0
%% body:
如图是一列简谐横波在某时刻的波形图，已知图中 b 位置的质点起振比 a 位置的质点晚 $0.5\Us$，b 和 c 之间的距离是 $5\Ucm$，则此列波的波长和频率应分别为\xzanswer{A}

\onepicture[w=5.2cm]{figs/12.png}

\fourchoices[answer=A]
{$5\Um$，$1\UHz$}
{$10\Um$，$2\UHz$}
{$5\Um$，$2\UHz$}
{$10\Um$，$1\UHz$}

%% answer: A
%% memo:
\memoanswer{
本题原卷及材料中未提供详解，待补充。
}

\item
%% number: 13
%% paperName: 2007年广东高考第 13 题 12 分
%% typeId: 4
%% type: 实验题
%% chapter: 电路
%% point: 多用表的使用
%% method: 无
%% score: 12
%% degree: 550
%% duplicateId: 0
%% body:
实验室新进了一批低电阻的电磁螺线管，已知螺线管使用的金属丝电阻率 $\rho=1.7\times10^{-8}\UOm$。课外活动小组的同学设计了一个试验来测算螺线管使用的金属丝长度。他们选择了多用电表、电流表、电压表、开关、滑动变阻器、螺旋测微器（千分尺）、导线和学生电源等。

\begin{enumerate}
	\item
	他们使用多用电表粗测金属丝的电阻，操作过程分以下三个步骤：（请填写第②步操作）

	①将红、黑表笔分别插入多用电表的“＋”、“－”插孔；选择电阻档“×1”；

	②\tkanswer[3.2cm]{将红、黑表笔短接，调整调零旋钮调零}；

	③把红黑表笔分别与螺线管金属丝的两端相接，多用表的示数如图 1（a）所示。

	\item
	根据多用电表示数，为了减少实验误差，并在实验中获得较大的电压调节范围，应从图 1（b）的 A、B、C、D 四个电路中选择 \tkanswer[1.2cm]{D} 电路来测量金属丝电阻；

	\item
	他们使用千分尺测量金属丝的直径，示数如图 2 所示，金属丝的直径为 \tkanswer[1.4cm]{0.260}$\Umm$；

	\item
	根据多用电表测得的金属丝电阻值，可估算出绕制这个螺线管所用金属丝的长度约为 \tkanswer[1.4cm]{12 或 13}$\Um$。（结果保留两位有效数字）

	\item
	他们正确连接电路，接通电源后，调节滑动变阻器，发现电流始终无示数。请设计一种方案，利用多用电表检查电路故障并写出判断依据。（只需写出简要步骤）

	\rule{12cm}{0.4pt}
\end{enumerate}

\threepicture[showlabel=false, align=right, hA=5.0cm, hB=5.0cm, hC=2.6cm]
{figs/13a.png}
{figs/13b.png}
{figs/13c.png}

%% answer: 
\jdanswer{
\begin{enumerate}
	\item
	将红、黑表笔短接，调整调零旋钮调零

	\item
	D

	\item
	$0.260\Umm$（$0.258\sim0.262\Umm$ 均给分）

	\item
	$12\Um$ 或 $13\Um$

	\item
	以下两种解答都正确：①使用多用电表的电压档位，接通电源，逐个测量各元件、导线上的电压，若电压等于电源电压，说明该元件或导线断路故障。②使用多用电表的电阻档位，断开电路或拆下元件、导线，逐个测量各元件、导线上的电阻，若电阻为无穷大，说明该元件或导线断路故障。

\end{enumerate}
}

%% memo:
\memoanswer{
本题原卷及材料中未提供详解，待补充。
}

\item
%% number: 14
%% paperName: 2007年广东高考第 14 题 8 分
%% typeId: 4
%% type: 实验题
%% chapter: 直线运动
%% point: 打点计时器公式
%% method: 无
%% score: 8
%% degree: 450
%% duplicateId: 0
%% body:
如图（a）所示，小车放在斜面上，车前端栓有不可伸长的细线，跨过固定在斜面边缘的小滑轮与重物相连，小车后面与打点计时器的纸带相连。起初小车停在靠近打点计时器的位置，重物到地面的距离小于小车到滑轮的的距离。启动打点计时器，释放重物，小车在重物的牵引下，由静止开始沿斜面向上运动，重物落地后，小车会继续向上运动一段距离。打点计时器使用的交流电频率为 $50\UHz$。图（b）中 a、b、c 是小车运动纸带上的三段，纸带运动方向如箭头所示。

\begin{enumerate}
	\item
	根据所提供纸带上的数据，计算打 c 段纸带时小车的加速度大小为 \tkanswer[1.6cm]{5.0}$\Umsq$。（结果保留两位有效数字）

	\item
	打 a 段纸带时，小车的加速度是 $2.5\Umsq$。请根据加速度的情况，判断小车运动的最大速度可能出现在 b 段纸带中的 \tkanswer[2.2cm]{$D_{4}D_{3}$ 区间内}。

	\item
	如果取重力加速度 $10\Umsq$，由纸带数据可推算出重物与小车的质量比为 \tkanswer[1.4cm]{$1:1$}。

\end{enumerate}

\twopicture[showlabel=false, align=right, wA=6.5cm, wB=11.0cm]
{figs/14a.png}
{figs/14b.png}

%% answer: 
\jdanswer{
\begin{enumerate}
	\item
	$5.0\Umsq$（结果是 $4.8\Umsq$ 的得 1 分）

	\item
	$D_{4}D_{3}$ 区间内

	\item
	$1:1$

\end{enumerate}
}

%% memo:
\memoanswer{
本题原卷及材料中未提供详解，待补充。
}

\item
%% number: 15
%% paperName: 2007年广东高考第 15 题 10 分
%% typeId: 4
%% type: 实验题
%% chapter: 磁场
%% point: 带电粒子在复合场中的运动
%% method: 无
%% score: 10
%% degree: 450
%% duplicateId: 0
%% body:
\begin{enumerate}
	\item
	放射性物质 $\ce{^{210}_{84}Po}$ 和 $\ce{^{60}_{27}Co}$ 的核衰变方程为：

	$\ce{^{210}_{84}Po -> ^{206}_{82}Pb + X_{1}}$，$\ce{^{60}_{27}Co -> ^{60}_{28}Ni + X_{2}}$。

	方程中的 $X_{1}$ 代表的是 \tkanswer[2.4cm]{$\ce{^{4}_{2}He}$（或 $\alpha$）}，$X_{2}$ 代表的是 \tkanswer[2.4cm]{$\ce{^{0}_{-1}e}$（或 $\beta$）}。

	\item
	如图 1（a）、（b）所示，铅盒内装有能释放 $\alpha$、$\beta$ 和 $\gamma$ 射线的放射性物质，在靠近铅盒的顶部加上电场 $E$ 或磁场 $B$，在图 1（a）、（b）中分别画出射线运动轨迹的示意图。（在所画的轨迹上须标明是 $\alpha$、$\beta$ 和 $\gamma$ 中的哪种射线）

	\item
	带电粒子的荷质比 $\frac{q}{m}$ 是一个重要的物理量。某中学物理兴趣小组设计了一个实验，探究电场和磁场对电子运动轨迹的影响，以求得电子的荷质比，实验装置如图 2 所示。

	①他们的主要实验步骤如下：

	A．首先在两极板 $M_{1}M_{2}$ 之间不加任何电场、磁场，开启阴极射线管电源，发射的电子从两极板中央通过，在荧幕的正中心处观察到一个亮点；

	B．在 $M_{1}M_{2}$ 两极板间加合适的电场：加极性如图 2 所示的电压，并逐步调节增大，使荧幕上的亮点逐渐向荧幕下方偏移，直到荧幕上恰好看不见亮点为止，记下此时外加电压为 $U$。请问本步骤目的是什么？

	C．保持步骤 B 中的电压 $U$ 不变，对 $M_{1}M_{2}$ 区域加一个大小、方向合适的磁场 $B$，使荧幕正中心重现亮点，试问外加磁场的方向如何？

	②根据上述实验步骤，同学们正确推算出电子的荷质比与外加电场、磁场及其他相关量的关系为 $\frac{q}{m}=\frac{U}{B^{2}d^{2}}$。一位同学说，这表明电子的荷质比将由外加电压决定，外加电压越大则电子的荷质比越大，你认为他的说法正确吗？为什么？

\end{enumerate}

\threepicture[numA={（a）}, numB={（b）}, numC={图2}, labelprefix={}, align=right, hA=4.0cm, hB=4.0cm, wC=9.5cm]
{figs/15a.png}
{figs/15b.png}
{figs/15c.png}

%% answer: 
\jdanswer{
\begin{enumerate}
	\item
	$X_{1}$ 代表的是 $\ce{^{4}_{2}He}$（或 $\alpha$），$X_{2}$ 代表的是 $\ce{^{0}_{-1}e}$（或 $\beta$）

	\item
	如答图所示。（曲率半径不作要求，每种射线可只画出一条轨迹）

	\item
	①B．使电子刚好落在正极板的近荧光屏端边缘，利用已知量表达 $\frac{q}{m}$；C．垂直电场方向向外（垂直纸面向外）。②说法不正确，电子的荷质比是电子的固有参数。

\end{enumerate}
}

%% memo:
\memoanswer{
本题原卷及材料中未提供详解，待补充。第（2）问的参考轨迹如答图。

\twopicture[showlabel=false, h=4.0cm]
{figs/15_an1.png}
{figs/15_an2.png}
}

\item
%% number: 16
%% paperName: 2007年广东高考第 16 题 12 分
%% typeId: 6
%% type: 计算题
%% chapter: 曲线运动
%% point: 天体运动
%% method: 无
%% score: 12
%% degree: 450
%% duplicateId: 0
%% body:
土星周围有许多大小不等的岩石颗粒，其绕土星的运动可视为圆周运动。其中有两个岩石颗粒 A 和 B 与土星中心距离分别为 $r_{A}=8.0\times10^{4}\Ukm$ 和 $r_{B}=1.2\times10^{5}\Ukm$。忽略所有岩石颗粒间的相互作用。（结果可用根式表示）

\begin{enumerate}
	\item
	求岩石颗粒 A 和 B 的线速度之比；
	\item
	求岩石颗粒 A 和 B 的周期之比；
	\item
	土星探测器上有一物体，在地球上重为 $10\UN$，推算出他在距土星中心 $3.2\times10^{5}\Ukm$ 处受到土星的引力为 $0.38\UN$。已知地球半径为 $6.4\times10^{3}\Ukm$，请估算土星质量是地球质量的多少倍？
\end{enumerate}

%% answer: 
\jdanswer{
\begin{enumerate}
	\item
	$\frac{\sqrt{6}}{2}$

	\item
	$\frac{2\sqrt{6}}{9}$

	\item
	$95$

\end{enumerate}
}

%% memo:
\memoanswer{
（1）设土星质量为 $M_{0}$，颗粒质量为 $m$，颗粒距土星中心距离为 $r$，线速度为 $v$，根据牛顿第二定律和万有引力定律：$\frac{GM_{0}m}{r^{2}}=\frac{mv^{2}}{r}$，解得 $v=\sqrt{\frac{GM_{0}}{r}}$。对于 A、B 两颗粒分别有 $v_{A}=\sqrt{\frac{GM_{0}}{r_{A}}}$ 和 $v_{B}=\sqrt{\frac{GM_{0}}{r_{B}}}$，得 $\frac{v_{A}}{v_{B}}=\frac{\sqrt{6}}{2}$。

（2）设颗粒绕土星作圆周运动的周期为 $T$，则 $T=\frac{2\pi r}{v}$。对于 A、B 两颗粒分别有 $T_{A}=\frac{2\pi r_{A}}{v_{A}}$ 和 $T_{B}=\frac{2\pi r_{B}}{v_{B}}$，得 $\frac{T_{A}}{T_{B}}=\frac{2\sqrt{6}}{9}$。

（3）设地球质量为 $M$，地球半径为 $r_{0}$，地球上物体的重力可视为万有引力，探测器上物体质量为 $m_{0}$，在地球表面重力为 $G_{0}$，距土星中心 $r_{0}'=3.2\times10^{5}\Ukm$ 处的引力为 $G_{0}'$，根据万有引力定律 $G_{0}=\frac{GMm_{0}}{r_{0}^{2}}$，$G_{0}'=\frac{GM_{0}m_{0}}{r_{0}'^{2}}$，解得 $\frac{M_{0}}{M}=95$。
}

\item
%% number: 17
%% paperName: 2007年广东高考第 17 题 16 分
%% typeId: 6
%% type: 计算题
%% chapter: 运动和力
%% point: 动量和能量
%% method: 无
%% score: 16
%% degree: 450
%% duplicateId: 0
%% body:
如图所示，在同一竖直面上，质量为 $2m$ 的小球 A 静止在光滑斜面的底部，斜面高度为 $H=2L$。小球受到弹簧的弹性力作用后，沿斜面向上运动。离开斜面后，达到最高点时与静止悬挂在此处的小球 B 发生弹性碰撞，碰撞后球 B 刚好能摆到与悬点 O 同一高度，球 A 沿水平方向抛射落在水平面 C 上的 P 点，O 点的投影 $O'$ 与 P 的距离为 $\frac{L}{2}$。已知球 B 质量为 $m$，悬绳长 $L$，视两球为质点，重力加速度为 $g$，不计空气阻力，求：

\begin{enumerate}
	\item
	球 B 在两球碰撞后一瞬间的速度大小；
	\item
	球 A 在两球碰撞前一瞬间的速度大小；
	\item
	弹簧的弹性力对球 A 所做的功。
\end{enumerate}

\onepicture[w=6.0cm, align=right]{figs/17.png}

%% answer: 
\jdanswer{
\begin{enumerate}
	\item
	$v_{B}=\sqrt{2gL}$

	\item
	$v_{0}=\frac{3}{4}\sqrt{2gL}$

	\item
	$W=\frac{57}{8}mgL$

\end{enumerate}
}

%% memo:
\memoanswer{
（1）碰撞后，根据机械能守恒定律，对 B 球有 $mgL=\frac{1}{2}mv_{B}^{2}$，解得 $v_{B}=\sqrt{2gL}$。

（2）A、B 球碰撞有 $2mv_{0}=2mv_{A}+mv_{B}$，$\frac{1}{2}\cdot2mv_{0}^{2}=\frac{1}{2}\cdot2mv_{A}^{2}+\frac{1}{2}\cdot mv_{B}^{2}$，解得 $v_{A}=\frac{1}{4}\sqrt{2gL}$，$v_{0}=\frac{3}{4}\sqrt{2gL}$。

（3）碰后 A 球做平抛运动，设平抛高度为 $y$，有 $\frac{L}{2}=v_{A}t$，$y=\frac{1}{2}gt^{2}$，解得 $y=L$。对 A 球应用动能定理得 $W-2mg(y+2L)=\frac{1}{2}\cdot2mv_{0}^{2}$，解得 $W=\frac{57}{8}mgL$。
}

\item
%% number: 18
%% paperName: 2007年广东高考第 18 题 17 分
%% typeId: 6
%% type: 计算题
%% chapter: 电磁感应
%% point: 电磁感应能量分析
%% method: 分情况讨论
%% score: 17
%% degree: 350
%% duplicateId: 0
%% body:
如图（a）所示，一端封闭的两条平行光滑导轨相距 $L$，距左端 $L$ 处的中间一段被弯成半径为 $H$ 的 $1/4$ 圆弧，导轨左右两段处于高度相差 $H$ 的水平面上。圆弧导轨所在区域无磁场，右段区域存在磁场 $B_{0}$，左段区域存在均匀分布但随时间线性变化的磁场 $B(t)$，如图（b）所示，两磁场方向均竖直向上。在圆弧顶端，放置一质量为 $m$ 的金属棒 ab，与导轨左段形成闭合回路，从金属棒下滑开始计时，经过时间 $t_{0}$ 滑到圆弧顶端。设金属棒在回路中的电阻为 $R$，导轨电阻不计，重力加速度为 $g$。

\begin{enumerate}
	\item
	问金属棒在圆弧内滑动时，回路中感应电流的大小和方向是否发生改变？为什么？
	\item
	求 0 到时间 $t_{0}$ 内，回路中感应电流产生的焦耳热量。
	\item
	探讨在金属棒滑到圆弧底端进入匀强磁场 $B_{0}$ 的一瞬间，回路中感应电流的大小和方向。
\end{enumerate}

\twopicture[numA={（a）}, numB={（b）}, labelprefix={}, align=right, wA=7.5cm, wB=5.0cm]
{figs/18a.png}
{figs/18b.png}

%% answer: 
\jdanswer{
\begin{enumerate}
	\item
	感应电流的大小和方向均不发生改变。因为金属棒滑到圆弧任意位置时，回路中磁通量的变化率相同。

	\item
	$Q=\frac{L^{4}B_{0}^{2}}{t_{0}R}$

	\item
	Ⅰ、当 $\sqrt{2gH}=\frac{L}{t_{0}}$ 时，$I=0$；Ⅱ、当 $\sqrt{2gH}>\frac{L}{t_{0}}$ 时，$I=\frac{B_{0}L}{R}\left(\sqrt{2gH}-\frac{L}{t_{0}}\right)$，方向为 b$\to$a；Ⅲ、当 $\sqrt{2gH}<\frac{L}{t_{0}}$ 时，$I=\frac{B_{0}L}{R}\left(\frac{L}{t_{0}}-\sqrt{2gH}\right)$，方向为 a$\to$b。

\end{enumerate}
}

%% memo:
\memoanswer{
（1）感应电流的大小和方向均不发生改变。因为金属棒滑到圆弧任意位置时，回路中磁通量的变化率相同。

（2）$0\sim t_{0}$ 时间内，设回路中感应电动势大小为 $E_{0}$，感应电流为 $I$，感应电流产生的焦耳热为 $Q$，由法拉第电磁感应定律 $E_{0}=\frac{\Delta\phi}{\Delta t}=L^{2}\frac{B_{0}}{t_{0}}$，根据闭合电路的欧姆定律 $I=\frac{E_{0}}{R}$，由焦耳定律有 $Q=I^{2}Rt=\frac{L^{4}B_{0}^{2}}{t_{0}R}$，解得 $Q=\frac{L^{4}B_{0}^{2}}{t_{0}R}$。

（3）设金属棒进入磁场 $B_{0}$ 一瞬间的速度为 $v$，金属棒在圆弧区域下滑的过程中机械能守恒 $mgH=\frac{1}{2}mv^{2}$。在很短的时间 $\Delta t$ 内，根据法拉第电磁感应定律，金属棒进入磁场 $B_{0}$ 区域瞬间的感应电动势为 $E$，则 $E=\frac{\Delta\phi}{\Delta t}$，$v=\frac{\Delta x}{\Delta t}$，$\Delta\phi=B_{0}L\Delta x+L^{2}\Delta B(t)$，由闭合电路欧姆定律得 $I=\frac{E}{R}$，解得感应电流 $I=\frac{B_{0}L}{R}\left(\sqrt{2gH}-\frac{L}{t_{0}}\right)$。根据上式讨论：Ⅰ、当 $\sqrt{2gH}=\frac{L}{t_{0}}$ 时，$I=0$；Ⅱ、当 $\sqrt{2gH}>\frac{L}{t_{0}}$ 时，$I=\frac{B_{0}L}{R}\left(\sqrt{2gH}-\frac{L}{t_{0}}\right)$，方向为 b$\to$a；Ⅲ、当 $\sqrt{2gH}<\frac{L}{t_{0}}$ 时，$I=\frac{B_{0}L}{R}\left(\frac{L}{t_{0}}-\sqrt{2gH}\right)$，方向为 a$\to$b。
}

\item
%% number: 19
%% paperName: 2007年广东高考第 19 题 17 分
%% typeId: 6
%% type: 计算题
%% chapter: 电场
%% point: 带电粒子的往复运动
%% method: 无
%% score: 17
%% degree: 350
%% duplicateId: 0
%% body:
如图所示，沿水平方向放置一条平直光滑槽，它垂直穿过开有小孔的两平行薄板，板相距 $3.5L$。槽内有两个质量均为 $m$ 的小球 A 和 B，球 A 带电量为 $+2q$，球 B 带电量为 $-3q$，两球由长为 $2L$ 的轻杆相连，组成一带电系统。最初 A 和 B 分别静止于左板的两侧，离板的距离均为 $L$。若视小球为质点，不计轻杆的质量，在两板间加上与槽平行向右的匀强电场 $E$ 后（设槽和轻杆由特殊绝缘材料制成，不影响电场的分布），求：

\begin{enumerate}
	\item
	球 B 刚进入电场时，带电系统的速度大小；
	\item
	带电系统从开始运动到速度第一次为零所需的时间及球 A 相对右板的位置。
\end{enumerate}

\onepicture[w=6.0cm, align=right]{figs/19.png}

%% answer: 
\jdanswer{
\begin{enumerate}
	\item
	$v_{1}=\sqrt{\frac{2qEL}{m}}$

	\item
	$t=\frac{7}{3}\sqrt{\frac{2mL}{qE}}$，$x=\frac{L}{6}$

\end{enumerate}
}

%% memo:
\memoanswer{
对带电系统进行分析，假设球 A 能达到右极板，电场力对系统做功为 $W_{1}$，有 $W_{1}=2qE\times2.5L+(-3qE\times1.5L)>0$，而且还能穿过小孔，离开右极板。假设球 B 能达到右极板，电场力对系统做功为 $W_{2}$，有 $W_{2}=2qE\times2.5L+(-3qE\times3.5L)<0$。综上所述，带电系统速度第一次为零时，球 A、B 应分别在右极板两侧。

（1）带电系统开始运动时，设加速度为 $a_{1}$，由牛顿第二定律 $a_{1}=\frac{2qE}{2m}=\frac{qE}{m}$。球 B 刚进入电场时，带电系统的速度为 $v_{1}$，有 $v_{1}^{2}=2a_{1}L$，求得 $v_{1}=\sqrt{\frac{2qEL}{m}}$。

（2）设球 B 从静止到刚进入电场的时间为 $t_{1}$，则 $t_{1}=\frac{v_{1}}{a_{1}}$，解得 $t_{1}=\sqrt{\frac{2mL}{qE}}$。球 B 进入电场后，带电系统的加速度为 $a_{2}$，由牛顿第二定律 $a_{2}=\frac{-3qE+2qE}{2m}=-\frac{qE}{2m}$。显然，带电系统做匀减速运动。设球 A 刚达到右极板时的速度为 $v_{2}$，减速所需时间为 $t_{2}$，则有 $v_{2}^{2}-v_{1}^{2}=2a_{2}\times1.5L$，$t_{2}=\frac{v_{2}-v_{1}}{a_{2}}$，求得 $v_{2}=\frac{1}{2}\sqrt{\frac{2qEL}{m}}$，$t_{2}=\sqrt{\frac{2mL}{qE}}$。球 A 离开电场后，带电系统继续做减速运动，设加速度为 $a_{3}$，再由牛顿第二定律 $a_{3}=\frac{-3qE}{2m}$。设球 A 从离开电场到静止所需的时间为 $t_{3}$，运动的位移为 $x$，则有 $t_{3}=\frac{0-v_{2}}{a_{3}}$，$-v_{2}^{2}=2a_{3}x$，求得 $t_{3}=\frac{1}{3}\sqrt{\frac{2mL}{qE}}$。可知，带电系统从静止到速度第一次为零所需的时间为 $t=t_{1}+t_{2}+t_{3}=\frac{7}{3}\sqrt{\frac{2mL}{qE}}$，球 A 相对右板的位置为 $x=\frac{L}{6}$。
}

\item
%% number: 20
%% paperName: 2007年广东高考第 20 题 18 分
%% typeId: 6
%% type: 计算题
%% chapter: 磁场
%% point: 带电粒子在磁场中的运动
%% method: 无
%% score: 18
%% degree: 350
%% duplicateId: 0
%% body:
如图是某装置的垂直截面图，虚线 $A_{1}A_{2}$ 是垂直截面与磁场区边界面的交线，匀强磁场分布在 $A_{1}A_{2}$ 的右侧区域，磁感应强度 $B=0.4\UT$，方向垂直纸面向外，$A_{1}A_{2}$ 与垂直截面上的水平线夹角为 $45\Udu$。在 $A_{1}A_{2}$ 左侧，固定的薄板和等大的挡板均水平放置，它们与垂直截面交线分别为 $S_{1}$、$S_{2}$，相距 $L=0.2\Um$。在薄板上 P 处开一小孔，P 与 $A_{1}A_{2}$ 线上点 D 的水平距离为 $L$。在小孔处装一个电子快门。起初快门开启，一旦有带正电微粒通过小孔，快门立即关闭，此后每隔 $T=3.0\times10^{-3}\Us$ 开启一次并瞬间关闭。从 $S_{1}S_{2}$ 之间的某一位置水平发射一速度为 $v_{0}$ 的带正电微粒，它经过磁场区域后入射到 P 处小孔。通过小孔的微粒与档板发生碰撞而反弹，反弹速度大小是碰前的 0.5 倍。

\begin{enumerate}
	\item
	经过一次反弹直接从小孔射出的微粒，其初速度 $v_{0}$ 应为多少？
	\item
	求上述微粒从最初水平射入磁场到第二次离开磁场的时间。（忽略微粒所受重力影响，碰撞过程无电荷转移。已知微粒的荷质比 $\frac{q}{m}=1.0\times10^{3}\UCkg$。只考虑纸面上带电微粒的运动）
\end{enumerate}

\onepicture[w=6.5cm, align=right]{figs/20.png}

%% answer: 
\jdanswer{
\begin{enumerate}
	\item
	$v_{0}=100\Ums$

	\item
	$t=2.8\times10^{-2}\Us$

\end{enumerate}
}

%% memo:
\memoanswer{
（1）如图 2 所示，设带正电微粒在 $S_{1}S_{2}$ 之间任意点 Q 以水平速度 $v_{0}$ 进入磁场，微粒受到的洛伦兹力为 $f$，在磁场中做圆周运动的半径为 $r$，有 $qv_{0}B=\frac{mv_{0}^{2}}{r}$，解得 $r=\frac{mv_{0}}{qB}$。欲使微粒能进入小孔，半径 $r$ 的取值范围为 $L<r<2L$，代入数据得 $80\Ums<v_{0}<160\Ums$。欲使进入小孔的微粒与挡板一次相碰返回后能通过小孔，还必须满足条件 $\frac{L}{v_{0}}+\frac{L}{0.5v_{0}}=nT$，其中 $n=1,2,3,\cdots$。可知，只有 $n=2$ 满足条件，即有 $v_{0}=100\Ums$。

（2）设微粒在磁场中做圆周运动的周期为 $T_{0}$，从水平进入磁场到第二次离开磁场的总时间为 $t$，设 $t_{1}$、$t_{4}$ 分别为带电微粒第一次、第二次在磁场中运动的时间，第一次离开磁场运动到挡板的时间为 $t_{2}$，碰撞后再返回磁场的时间为 $t_{3}$，运动轨迹如答图 2 所示，则有 $T_{0}=\frac{2\pi r}{v_{0}}$，$t_{1}=\frac{3}{4}T_{0}$，$t_{2}=\frac{2L}{v_{0}}$，$t_{3}=\frac{2L}{0.5v_{0}}$，$t_{4}=\frac{1}{4}T_{0}$，$t=t_{1}+t_{2}+t_{3}+t_{4}=2.8\times10^{-2}\Us$。

\onepicture[w=6.0cm]{figs/20_an.png}
}

\end{enumerate}

\end{document}
